% =====================================================================
%  THÈME 5 — Rendement et bilan complet
%  Niveau MOYEN — Corrigé
% =====================================================================
\documentclass[11pt,a4paper]{article}
\usepackage[utf8]{inputenc}
\usepackage[T1]{fontenc}
\usepackage{lmodern}
\usepackage[french]{babel}
\usepackage{amsmath,amssymb,mathtools}
\usepackage{icomma}
\usepackage[version=4]{mhchem}
\usepackage{siunitx}
\sisetup{output-decimal-marker={,},per-mode=symbol}
\usepackage{xcolor}
\usepackage{tikz}
\usepackage[most]{tcolorbox}
\usepackage{enumitem}
\usepackage{array}
\usepackage{fancyhdr}
\usepackage[a4paper,margin=2.1cm,headheight=26pt,headsep=8pt]{geometry}

\definecolor{primary}{RGB}{142,93,168}
\definecolor{primarydark}{RGB}{82,45,108}
\definecolor{corrcol}{RGB}{176,110,20}
\newcommand{\nq}[1]{n_{\text{#1}}}

\pagestyle{fancy}\fancyhf{}
\fancyhead[L]{\small\textcolor{primarydark}{\textbf{Fiche 5} — Thème 5 : Rendement}}
\fancyhead[R]{\small\textcolor{corrcol}{\textsc{Corrigé — Moyen}}}
\fancyfoot[C]{\small\thepage}
\renewcommand{\headrulewidth}{0.8pt}
\renewcommand{\headrule}{\hbox to\headwidth{\color{primary}\leaders\hrule height \headrulewidth\hfill}}

\newtcolorbox[auto counter]{cor}[1][]{enhanced,breakable,boxrule=1pt,arc=3pt,
  left=3mm,right=3mm,top=2.4mm,bottom=2.4mm,colback=corrcol!6,colframe=corrcol,
  fonttitle=\bfseries,coltitle=white,
  title={Corrigé~\thetcbcounter\ \ifx\relax#1\relax\relax\else\textmd{--- \itshape #1}\fi}}

\newcommand{\rep}[1]{\par\smallskip\noindent\textbf{\color{corrcol}$\Rightarrow$ #1}}

\setlength{\parindent}{0pt}
\setlength{\parskip}{3pt}

\begin{document}

\begin{center}
{\LARGE\bfseries\color{primarydark}Thème 5 — Rendement}\\[1pt]
{\large\color{corrcol}\textbf{Corrigé — Niveau Moyen}}\\
{\color{primary}\rule{0.6\linewidth}{1pt}}
\end{center}

\begin{cor}[fermentation du glucose]
\textbf{a)} Ratio \ce{C6H12O6}:\ce{C2H5OH}$=1:2$ :
$\nq{éthanol,théo}=2\times0{,}50=1{,}0$~mol, soit $m_{\text{théo}}=1{,}0\times46=\SI{46}{\gram}$.\\
\textbf{b)} $R=\dfrac{m_{\text{réel}}}{m_{\text{théo}}}\times100=\dfrac{11{,}5}{46}\times100=\boxed{\SI{25}{\percent}}$.
\rep{Ou en moles : $n_{\text{réel}}=11{,}5/46=0{,}25$~mol, $R=0{,}25/1{,}0=25\,\%$.}
\end{cor}

\begin{cor}[aluminothermie du fer]
\textbf{a)} $\nq{FeO}=\dfrac{144}{72}=2$~mol ; $\nq{Al}=\dfrac{27}{27}=1$~mol. Rapports :
$\frac{2}{3}\approx0{,}667$ et $\frac{1}{2}=0{,}5$. \boxed{\text{\ce{Al} est limitant}}.\\
\textbf{b)} $\nq{Fe,théo}=\nq{Al}\times\dfrac{3}{2}=1\times1{,}5=1{,}5$~mol, soit
$m_{\text{théo}}=1{,}5\times56=\SI{84}{\gram}$.\\
\textbf{c)} $R=\dfrac{42}{84}\times100=\boxed{\SI{50}{\percent}}$.
\end{cor}

\begin{cor}[synthèse du carbure de silicium]
\textbf{a)} $\nq{SiO2}=\dfrac{120000}{60}=2000$~mol ; $\nq{C}=\dfrac{36000}{12}=3000$~mol. Rapports :
$\frac{2000}{1}=2000$ et $\frac{3000}{3}=1000$. \boxed{\text{le carbone \ce{C} est limitant}}.\\
\textbf{b)} $\nq{SiC,théo}=\nq{C}\times\dfrac{1}{3}=1000$~mol, soit
$m_{\text{théo}}=1000\times40=\SI{40000}{\gram}=\SI{40}{\kilo\gram}$.\\
\textbf{c)} $R=\dfrac{32}{40}\times100=\boxed{\SI{80}{\percent}}$.
\end{cor}

\begin{cor}[préparation du carbonate de sodium]
\textbf{a)} $\nq{NaCl}=\dfrac{29{,}25}{58{,}5}=0{,}50$~mol ; $\nq{CaCO3}=\dfrac{45{,}0}{100}=0{,}45$~mol.
Rapports : $\frac{0{,}50}{2}=0{,}25$ et $\frac{0{,}45}{1}=0{,}45$. \boxed{\text{\ce{NaCl} est limitant}}.\\
\textbf{b)} $\nq{Na2CO3,théo}=\nq{NaCl}\times\dfrac{1}{2}=0{,}25$~mol, soit
$m_{\text{théo}}=0{,}25\times106=\SI{26,5}{\gram}$.
\[
R=\frac{13{,}25}{26{,}5}\times100=\boxed{\SI{50}{\percent}}.
\]
\end{cor}

\begin{cor}[combustion du sulfure d'hydrogène avec pertes]
\textbf{a)} $\nq{H2S}=\dfrac{10{,}2}{34}=0{,}30$~mol ; $\nq{O2}=\dfrac{9{,}6}{32}=0{,}30$~mol. Rapports :
$\frac{0{,}30}{2}=0{,}15$ et $\frac{0{,}30}{3}=0{,}10$. \boxed{\text{\ce{O2} est limitant}}.\\
\textbf{b)} Théorique : $\nq{SO2,théo}=\nq{O2}\times\dfrac{2}{3}=0{,}30\times\dfrac{2}{3}=0{,}20$~mol.
Avec $R=\SI{50}{\percent}$ : $\nq{SO2,réel}=\dfrac{50}{100}\times0{,}20=0{,}10$~mol, soit
$m=0{,}10\times64=\boxed{\SI{6,4}{\gram}}$.
\end{cor}

\end{document}
